\documentclass[12pt,letterpaper]{article}

\usepackage[margin=1in]{geometry}
\usepackage{amsmath,amssymb}
\usepackage{fancyhdr}

\setlength{\parindent}{0pt}
\setlength{\parskip}{0.8em}
\setlength{\headheight}{15pt}

\pagestyle{fancy}
\fancyhf{}
\fancyhead[L]{FIT3155 SAMPLE EXAM PAPER}
\fancyfoot[C]{Page \thepage{} of 22}
\renewcommand{\headrulewidth}{0pt}
\renewcommand{\footrulewidth}{0pt}

\newcommand{\marksbox}[1]{%
  \vfill
  \begin{flushright}
    \fbox{\parbox[c][1.1cm][c]{1.3cm}{\centering #1}}
  \end{flushright}
}

\newcommand{\questiontitle}[2]{%
  \textbf{Question #1}\hfill\textbf{(#2 marks)}
}

\newcommand{\workingpage}{%
  \newpage
  \begin{center}
    (Page intentionally left blank for working space)
  \end{center}
}

\begin{document}

\begin{titlepage}
  \thispagestyle{empty}
  \vspace*{0.18\textheight}
  {\Large FIT3155 \hfill SAMPLE EXAM PAPER\par}
  \vspace{2.2cm}
  \begin{center}
    {\LARGE\bfseries FIT3155 SAMPLE EXAM PAPER\par}
    \vspace{0.8cm}
    {\large Generated Sample Exam 05\par}
  \end{center}
\end{titlepage}

\setcounter{page}{2}

\section*{INSTRUCTIONS}

\begin{itemize}
  \item This exam is worth 60 marks.
\end{itemize}

\begin{center}
  \textbf{General exam technique}
\end{center}

\begin{itemize}
  \item Some students lose marks by not attempting all questions.

  \item Suppose you are likely to get 3/6 on a question after spending say 10 minutes on it.
  Spending another 10 minutes on the same question may gain you at most 3 additional
  marks. In contrast, spending those 10 minutes on a new question could earn you another
  6 marks. The key point is: do not get stuck on one question.

  \item A good strategy is to first answer the questions you are most confident about before
  attempting the others. In other words, answering questions strictly serially during an
  exam can be suboptimal.

  \item Do not waste time answering an unrelated question in the hope of gaining partial marks.
  Marking will be based on the question that was actually asked.

  \item Where necessary, justify your answers with clear and concise explanations, without being
  overly wordy.
\end{itemize}

\newpage

\questiontitle{1}{6}

Let $txt[1 \ldots n]$ be a text and $pat[1 \ldots m]$ a pattern, both drawn
from an alphabet $\Sigma$, with $m \leq n$.

Write pseudocode for an efficient algorithm that outputs \emph{every} starting
position $j$ in $txt$ at which $pat$ occurs exactly, i.e.\ every $j$ such that
$txt[j \ldots j + m - 1] = pat[1 \ldots m]$. State the worst-case time
complexity of your algorithm and briefly justify it.

(For this problem, you may assume that \texttt{zalg} is available as a built-in
function that takes a string as input and returns the array of its Z-values.
Do not write pseudocode for \texttt{zalg}. \emph{Hint:} consider the Z-values
of a single carefully chosen concatenation, using a separator symbol that
appears in neither $pat$ nor $txt$.)

\marksbox{6}

\workingpage
\newpage

\questiontitle{2}{6}

In the Boyer--Moore algorithm for exact matching of the pattern
$pat[1 \ldots m]$, the good-suffix array $gs[1 \ldots m+1]$ is used to decide a
safe shift after a mismatch. The intended meaning is the one used in the
shifting rule: when a right-to-left scan mismatches at position $k$ (so the
suffix $pat[k+1 \ldots m]$ has just been matched), the pattern may be shifted
$m - gs[k+1]$ positions to the right, where $gs[k+1]$ is the right-end position
in $pat$ of the rightmost \emph{internal} copy of the matched suffix
$pat[k+1 \ldots m]$ whose immediately preceding character differs from
$pat[k]$ (and $gs[k+1] = 0$ if no such copy exists).

Write pseudocode that, given $pat[1 \ldots m]$, constructs the entire array
$gs[1 \ldots m+1]$. You may assume that the Z-suffix values
$Zsuf[1 \ldots m]$ of $pat$ are available, where $Zsuf_i$ is the length of the
longest substring of $pat$ \emph{ending} at position $i$ that is identical to a
suffix of $pat$ (these are simply the Z-values of the reversed pattern, read
back to front). State the worst-case time and space complexity of your
construction.

\marksbox{6}

\workingpage
\newpage

\questiontitle{3}{6}

Let $S[1 \ldots n-1]\$$ be a string terminated by a unique symbol \$ that is
lexicographically smaller than every other character, and let
$L[1 \ldots n]$ be its Burrows--Wheeler Transform (the last column of the
sorted matrix of cyclic rotations of $S$).

You are given only the transform
\[
  L = \texttt{rttaar\$}.
\]

Using the last-to-first (LF) mapping, invert this BWT to recover the original
string $S$. Show your working clearly: state the first column $F$ that you use,
explain how each successive character of $S$ is recovered, and present the
final reconstructed string.

\marksbox{6}

\workingpage
\newpage

\questiontitle{4}{6=3+3}

These questions relate to the suffix tree of a string $S[1 \ldots n]\$$, built
by Ukkonen's linear-time algorithm, where \$ is a unique terminal symbol.

\begin{enumerate}
  \item[(i)] Prove that the completed suffix tree has exactly $n$ leaves and at
  most $n - 1$ internal nodes (not counting the root). Your argument should
  make clear why every internal node has at least two children, and why the
  terminal symbol \$ guarantees there is one leaf per suffix.

  \item[(ii)] Define what a \emph{suffix link} is. Then explain how suffix
  links are used to move from one suffix extension to the next \emph{within a
  single phase} of Ukkonen's algorithm, and why their use (together with the
  skip/count trick) is essential to achieving the overall $O(n)$ running time.
\end{enumerate}

\marksbox{6}

\workingpage
\newpage

\questiontitle{5}{6=3+3}

\begin{enumerate}
  \item[(i)] Work out the Elias-Omega binary codeword for the positive integer
  $42$. Show each length component of your codeword and the order in which the
  components are concatenated.

  \item[(ii)] Encode the string
  \[
    S = \texttt{ababab}
  \]
  using the LZ77 method, producing the sequence of tuples
  $[\textit{offset}, \textit{length}, \textit{next character}]$. Assume an
  unbounded search window and unbounded lookahead, and allow a copy whose
  length exceeds its offset (overlapping copy). Show how each tuple is formed.
\end{enumerate}

\marksbox{6}

\workingpage
\newpage

\questiontitle{6}{6=3+3}

These questions concern the number-theoretic fact at the heart of the
Miller--Rabin primality test.

\begin{enumerate}
  \item[(i)] Prove that if $p$ is prime, then the only solutions of
  \[
    x^{2} \equiv 1 \pmod{p}
  \]
  are $x \equiv 1 \pmod{p}$ and $x \equiv p - 1 \pmod{p}$. (In other words, the
  only square roots of $1$ modulo a prime are the \emph{trivial} ones,
  $\pm 1$.)

  \item[(ii)] Write $n - 1 = 2^{s} t$ with $t$ odd. Using part (i), explain why
  the Miller--Rabin test, on examining the sequence
  $a^{t},\, a^{2t},\, a^{4t},\, \ldots,\, a^{2^{s} t} \pmod{n}$, can correctly
  declare $n$ composite. In particular, explain what makes a base $a$ a
  \emph{witness} to the compositeness of $n$.
\end{enumerate}

\marksbox{6}

\workingpage
\newpage

\questiontitle{7}{6}

Starting from an initially empty Fibonacci heap, the following keys are
inserted, one at a time, in this order:
\[
  7,\; 3,\; 8,\; 1,\; 5,\; 9,\; 2,\; 6.
\]
A single \textsc{Extract-Min} operation is then performed.

Use the following conventions: \textsc{Insert} lazily adds a new single-node
tree to the root list; after the minimum is removed, the consolidation step
processes the root list and repeatedly links two roots of \emph{equal} degree
by making the root with the larger key a child of the root with the smaller
key.

\begin{enumerate}
  \item[(i)] Draw the heap (i.e.\ the root list and the trees it contains) at
  the moment just \emph{before} \textsc{Extract-Min} is performed, and indicate
  \texttt{H.min}.

  \item[(ii)] Perform the \textsc{Extract-Min}. Draw the resulting heap
  \emph{after} consolidation, clearly showing every tree, the parent--child
  relationships, and the new \texttt{H.min}.
\end{enumerate}

\marksbox{6}

\workingpage
\newpage

\questiontitle{8}{6}

A first-in-first-out queue is implemented using two stacks, \textsc{In} and
\textsc{Out}, each supporting \textsc{Push} and \textsc{Pop} at unit cost per
element:

\begin{itemize}
  \item \textsc{Enqueue}$(x)$: push $x$ onto \textsc{In}.
  \item \textsc{Dequeue}$()$: if \textsc{Out} is empty, repeatedly pop every
  element off \textsc{In} and push it onto \textsc{Out}; then pop and return the
  top of \textsc{Out}.
\end{itemize}

Both stacks are initially empty. Using the \emph{accounting} (banker's) method
of amortized analysis, show that the amortized cost of both \textsc{Enqueue}
and \textsc{Dequeue} is $O(1)$. State the charge you assign to each operation,
state the credit invariant your scheme maintains, and verify that the total
stored credit can never become negative over any sequence of operations
starting from two empty stacks.

\marksbox{6}

\workingpage
\newpage

\questiontitle{9}{6=3+3}

Consider a B-tree with minimum degree $t$, in which every node stores between
$t - 1$ and $2t - 1$ keys (the root may hold fewer). Insertion uses the
top-down splitting strategy: whenever a full child (one holding $2t - 1$ keys)
is encountered on the way down, it is split before being entered.

\begin{enumerate}
  \item[(i)] Write pseudocode for the operation $\textsc{Split-Child}(x, i)$,
  which splits the full child $y$ (the $i$-th child of node $x$) into two nodes
  and promotes its median key into $x$. State the worst-case time complexity of
  $\textsc{Split-Child}$ in terms of $t$.

  \item[(ii)] Explain why $\textsc{Split-Child}$ preserves all the B-tree
  invariants (the bounds on the number of keys per node, and the requirement
  that all leaves remain at the same depth), and explain why, under this
  insertion strategy, the height of the B-tree can increase only when the root
  itself is split.
\end{enumerate}

\marksbox{6}

\workingpage
\newpage

\questiontitle{10}{6=3+3}

\begin{enumerate}
  \item[(i)] Solve the following linear program using the tableau simplex
  algorithm. Show the tableau after each pivot, and state the optimal values of
  $x$, $y$, and $z$.
  \[
  \begin{array}{rl}
  \text{Maximize} & z = 3x + 4y \\
  \text{Subject to} & 2x + y \leq 10, \\
                    & x + 3y \leq 15, \\
                    & x \geq 0,\quad y \geq 0.
  \end{array}
  \]

  \item[(ii)] Explain how, while running the tableau simplex method, one can
  recognise from a tableau that the linear program is \emph{unbounded} (i.e.\
  the objective can be increased without limit). Describe precisely the
  condition on a chosen entering column that signals unboundedness, and state
  what the algorithm should report in that case.
\end{enumerate}

\marksbox{6}

\workingpage

\vfill
\begin{center}
  End of Exam
\end{center}

\end{document}
